\chapter{Divisors and Line Bundles}\label{divisors}
\chapterstatus{complete}

\section{Introduction}\label{divisors:section-introduction}

Divisors are the codimension-one cycles. They are the first case of the algebraic cycles whose cohomology classes the Hodge conjecture is
about, and the case in which the conjecture is a theorem (the Lefschetz $(1,1)$-theorem, Chapter~\ref{lefschetz-11}). This chapter treats
Weil divisors and the class group, Cartier divisors and their identification with line bundles, linear systems and morphisms to projective
space, ample line bundles, intersection numbers on surfaces, and the N\'eron--Severi group. References: \cite[Sections II.6, II.7 and V.1]{hartshorne-ag},
\cite[Chapter 1]{lazarsfeld-positivity-1}, \cite[Chapter I]{beauville-surfaces}.

\noindent\textbf{Prerequisites.} Chapter~\ref{coherent-cohomology} (Cohomology of Coherent Sheaves).

Throughout, schemes are Noetherian and separated. $K(X)$ is the function field of an integral scheme and $\mathcal{K}_X$ the constant sheaf $K(X)$.

\section{Weil divisors}\label{divisors:section-weil}

\begin{definition}\label{divisors:definition-weil}
Let $X$ be an integral scheme that is \emph{regular in codimension one}: every local ring $\mathcal{O}_{X,x}$ of dimension one is regular, hence a discrete valuation ring
(Theorem~\ref{commutative-algebra:theorem-dvr}). A \emph{prime divisor} is a closed integral subscheme $Y$ of codimension one; its generic point $\eta_Y$ has local
ring a DVR with valuation $v_Y$ on $K(X)$. A \emph{Weil divisor} is an element $D = \sum n_YY$ of the free abelian group $\mathrm{Div}(X)$ on the prime divisors; it is
\emph{effective}, $D \geq 0$, if all $n_Y \geq 0$. For $f \in K(X)^*$ put $(f) = \sum_Yv_Y(f)Y$, a \emph{principal divisor}. Divisors $D, D'$ are \emph{linearly
equivalent}, $D \sim D'$, if $D - D'$ is principal, and the \emph{class group} is $\mathrm{Cl}(X) = \mathrm{Div}(X)/\{\text{principal divisors}\}$.
\end{definition}

\begin{lemma}\label{divisors:lemma-finite-sum}
For $f \in K(X)^*$, $v_Y(f) = 0$ for all but finitely many $Y$, so $(f)$ is a divisor, and $f \mapsto (f)$ is a homomorphism $K(X)^* \to \mathrm{Div}(X)$.
\end{lemma}

\begin{proof}
$f$ is a regular function with regular inverse on some nonempty open $U$. The complement $X \setminus U$ is a proper closed subset of the Noetherian space $X$, so
it contains only finitely many prime divisors (they are among its finitely many irreducible components). For $Y$ meeting $U$, $f$ is a unit at $\eta_Y$, so
$v_Y(f) = 0$. Additivity is $v_Y(fg) = v_Y(f) + v_Y(g)$.
\end{proof}

\begin{proposition}\label{divisors:proposition-class-group}
\begin{enumerate}
\item Let $A$ be a Noetherian domain. If $A$ is a unique factorisation domain, then $\mathrm{Cl}(\Spec A) = 0$. In particular $\mathrm{Cl}(\Aa^n_k) = 0$.
\item If $Z \subseteq X$ is closed and $U = X \setminus Z$, there is an exact sequence $\bigoplus\ZZ Z_i \to \mathrm{Cl}(X) \to \mathrm{Cl}(U) \to 0$, the sum running over the
irreducible components $Z_i$ of $Z$ of codimension one.
\item For $X = \PP^n_k$, the degree $\deg\sum n_YY = \sum n_Y\deg Y$, where $\deg Y$ is the degree of a homogeneous generator of the ideal of $Y$, induces
$\mathrm{Cl}(\PP^n_k) \cong \ZZ$, generated by a hyperplane $H$.
\end{enumerate}
\end{proposition}

\begin{proof}
(1) A prime divisor of $\Spec A$ corresponds to a prime $\mathfrak{p}$ of height one. It contains an irreducible factor $f$ of any of its nonzero elements, and $(f)$ is a
nonzero prime in a UFD, so $\mathfrak{p} = (f)$. Then $Y = (f)$ is principal, because $v_{(f)}(f) = 1$ and $v_{(g)}(f) = 0$ for irreducibles $g$ not associate to $f$.
(2) $Y \mapsto Y \cap U$ gives a surjection $\mathrm{Div}(X) \to \mathrm{Div}(U)$, since a prime divisor of $U$ is the trace of its closure. Its kernel is generated by the
$Z_i$, and principal divisors map to principal divisors, $(f)|_U = (f|_U)$. (3) A prime divisor of $\PP^n$ is $V_+(g)$ with $g$ irreducible homogeneous, by
Theorem~\ref{varieties:theorem-hypersurface-sections} applied to the cone. For $f = g/h$ with $g$, $h$ homogeneous of the same degree, $(f)$ has degree
$\deg g - \deg h = 0$, so $\deg$ is defined on $\mathrm{Cl}$. It is onto, as $\deg H = 1$. If $D = \sum n_iY_i$ with $Y_i = V_+(g_i)$ has degree $0$, then
$f = \prod_ig_i^{n_i}$ is a quotient of homogeneous polynomials of the same degree, hence an element of $K(\PP^n)^*$, and $(f) = D$ because $v_{Y_j}(g_i) = \delta_{ij}$.
So $\deg$ is injective.
\end{proof}

\begin{example}\label{divisors:example-quadric-cone}
Let $k$ be a field of characteristic $\neq 2$ and $X = \Spec k[x, y, z]/(xy - z^2)$, the quadric cone. It is normal, and the line $Y = \{y = z = 0\}$ is a prime divisor with
$2Y = (y)$, since $y$ vanishes to order $2$ along $Y$ (on $Y$, $x$ is a unit near the generic point and $y = z^2/x$). One shows $\mathrm{Cl}(X) \cong \ZZ/2\ZZ$,
generated by $Y$ \cite[Example II.6.5.2]{hartshorne-ag}. So $Y$ is not locally principal at the vertex, although $2Y$ is.
\end{example}

\section{Cartier divisors and invertible sheaves}\label{divisors:section-cartier}

\begin{definition}\label{divisors:definition-cartier}
Let $X$ be an integral scheme. A \emph{Cartier divisor} is a global section of $\mathcal{K}_X^*/\mathcal{O}_X^*$. It is given by an open cover $(U_i)$ and
$f_i \in K(X)^*$ with $f_i/f_j \in \mathcal{O}^*(U_i \cap U_j)$; the $f_i$ are \emph{local equations}. It is \emph{principal} if it is the image of some
$f \in K(X)^*$, and \emph{effective} if all $f_i \in \mathcal{O}(U_i)$. The group of Cartier divisors modulo principal ones is $\mathrm{CaCl}(X)$. For a Cartier divisor
$D$, $\mathcal{O}(D) \subseteq \mathcal{K}_X$ is the subsheaf generated by $f_i^{-1}$ on $U_i$.
\end{definition}

\begin{proposition}\label{divisors:proposition-cartier-picard}
Let $X$ be an integral scheme.
\begin{enumerate}
\item $\mathcal{O}(D)$ is invertible, $\mathcal{O}(D - D') \cong \mathcal{O}(D) \otimes \mathcal{O}(D')^{-1}$, and $D \mapsto \mathcal{O}(D)$ induces an isomorphism
$\mathrm{CaCl}(X) \cong \Pic(X)$.
\item If $X$ is moreover regular in codimension one and \emph{locally factorial}, meaning that all local rings are unique factorisation domains, for instance if $X$ is
regular (Remark~\ref{commutative-algebra:remark-regular-domain}), then Weil and Cartier divisors coincide, and $\mathrm{Cl}(X) \cong \mathrm{CaCl}(X) \cong \Pic(X)$.
\end{enumerate}
\end{proposition}

\begin{proof}
(1) On $U_i$, $\mathcal{O}(D)$ is free with generator $f_i^{-1}$, so it is invertible, and multiplication of local equations corresponds to tensor product of
sub-$\mathcal{O}_X$-modules of $\mathcal{K}_X$. $\mathcal{O}(D) \cong \mathcal{O}_X$ if and only if there is a global generator $f^{-1}$, if and only if $D$ is principal, so the map
is injective on classes. It is onto: an invertible sheaf $\mathcal{L}$ on an integral scheme embeds in $\mathcal{L} \otimes \mathcal{K}_X \cong \mathcal{K}_X$ (a constant sheaf, as
$X$ is irreducible), and a local generator $e_i$ of $\mathcal{L}$ becomes some $f_i^{-1} \in K(X)^*$, so $\mathcal{L} \cong \mathcal{O}(D)$ with local equations $f_i$.
(2) A Cartier divisor gives the Weil divisor $\sum_Yv_Y(f_Y)Y$, where $f_Y$ is a local equation near $\eta_Y$. Conversely a Weil divisor is locally principal:
near $x$, each prime divisor through $x$ corresponds to a height-one prime of the UFD $\mathcal{O}_{X,x}$, which is principal, generated by some $g$. So $D$ has
the local equation $\prod g^{n_Y}$ on a neighbourhood of $x$. These constructions are inverse, and principal divisors correspond.
\end{proof}

\begin{example}\label{divisors:example-picard-projective}
$\Pic(\PP^n_k) \cong \ZZ$, generated by $\mathcal{O}(1) = \mathcal{O}(H)$, by Propositions~\ref{divisors:proposition-class-group}(3) and~\ref{divisors:proposition-cartier-picard}.
For the quadric cone of Example~\ref{divisors:example-quadric-cone}, $\Pic(X) = 0$ while $\mathrm{Cl}(X) = \ZZ/2$: the line $Y$ is a Weil divisor that is not Cartier.
\end{example}

\section{Linear systems and morphisms to projective space}\label{divisors:section-linear-systems}

\begin{proposition}\label{divisors:proposition-morphisms-projective}
Let $X$ be a scheme over a ring $A$. Morphisms $\varphi \colon X \to \PP^n_A$ over $A$ correspond bijectively to invertible sheaves $\mathcal{L}$ with sections
$s_0, \ldots, s_n \in \Gamma(X, \mathcal{L})$ generating $\mathcal{L}$, up to isomorphism of such data. The correspondence sends $\varphi$ to $\varphi^*\mathcal{O}(1)$ and
$s_i = \varphi^*x_i$.
\end{proposition}

\begin{proof}
Given $(\mathcal{L}, s_i)$, the open sets $X_{s_i}$ cover $X$ because the $s_i$ generate. On $X_{s_i}$ the quotients $s_j/s_i$ are regular functions, and the
ring map $A[x_0/x_i, \ldots, x_n/x_i] \to \Gamma(X_{s_i}, \mathcal{O})$, $x_j/x_i \mapsto s_j/s_i$, defines $X_{s_i} \to U_i$
(Theorem~\ref{schemes:theorem-morphisms-to-affine}). These glue, since the quotients satisfy the relations of the gluing of $\PP^n$. Then
$\varphi^*\mathcal{O}(1) \cong \mathcal{L}$, the generator $x_i$ of $\mathcal{O}(1)$ on $U_i$ being sent to $s_i$. Conversely $\varphi^*x_i$ generate $\varphi^*\mathcal{O}(1)$, and
$\varphi$ is recovered on each $\varphi^{-1}(U_i) = X_{\varphi^*x_i}$ from the quotients.
\end{proof}

\begin{definition}\label{divisors:definition-linear-system}
Let $X$ be a projective variety over an algebraically closed field $k$ and $D$ a Cartier divisor. The \emph{complete linear system} $|D|$ is the set of effective divisors
linearly equivalent to $D$. A \emph{linear system} is a subset corresponding to a linear subspace of $\Gamma(X, \mathcal{O}(D))$ as below, and its \emph{base points}
are the common points of its members.
\end{definition}

\begin{proposition}\label{divisors:proposition-linear-system}
In the situation of Definition~\ref{divisors:definition-linear-system}, sending a nonzero $s \in \Gamma(X, \mathcal{O}(D))$ to its divisor of zeros $(s)_0$ gives a bijection
$(\Gamma(X, \mathcal{O}(D)) \setminus 0)/k^* \cong |D|$. So $|D|$ is a projective space of dimension $h^0(\mathcal{O}(D)) - 1$, finite by
Corollary~\ref{coherent-cohomology:corollary-finite-dimensional}.
\end{proposition}

\begin{proof}
A section $s$ of $\mathcal{O}(D) \subseteq \mathcal{K}$ is an element $f \in K(X)$ with $f_if$ regular on each $U_i$, and $(s)_0$ is the effective divisor with local equations
$f_if$, which equals $D + (f)$. Conversely, if $D' = D + (f) \geq 0$, then $f \in \Gamma(X, \mathcal{O}(D))$. Two sections give the same divisor if and only if their quotient
is a unit on $X$, that is, an element of $\Gamma(X, \mathcal{O}_X)^* = k^*$ (Corollary~\ref{coherent-cohomology:corollary-finite-dimensional}).
\end{proof}

\begin{proposition}\label{divisors:proposition-closed-immersion-criterion}
Let $X$ be a projective scheme over an algebraically closed field $k$ and $\varphi \colon X \to \PP^n_k$ the morphism given by sections $s_0, \ldots, s_n$ generating
$\mathcal{L}$, spanning $V \subseteq \Gamma(X, \mathcal{L})$. Then $\varphi$ is a closed immersion if and only if $V$ separates points (for closed $P \neq Q$ there is $s \in V$
vanishing at $P$ but not at $Q$) and separates tangent vectors (for every closed $P$, the sections vanishing at $P$ span $\mathfrak{m}_P\mathcal{L}_P/\mathfrak{m}_P^2\mathcal{L}_P$).
\end{proposition}

This is \cite[Proposition II.7.3]{hartshorne-ag}; it is the algebraic version of Lemma~\ref{kodaira:lemma-embedding-criterion}. The proof uses that a proper
morphism that is injective on closed points and on tangent spaces is a closed immersion, by Nakayama's lemma applied to $\varphi_*\mathcal{O}_X$.

\section{Ample line bundles}\label{divisors:section-ample}

\begin{theorem}\label{divisors:theorem-ample-very-ample}
Let $X$ be of finite type over a Noetherian ring $A$ and $\mathcal{L}$ invertible. Then $\mathcal{L}$ is ample (Definition~\ref{coherent-sheaves:definition-very-ample}) if and only
if $\mathcal{L}^m$ is very ample over $A$ for some $m > 0$.
\end{theorem}

\begin{proof}
If $\mathcal{L}^m$ is very ample, $X$ is an open subscheme of a projective scheme $\bar X$ with $\mathcal{L}^m = \mathcal{O}(1)|_X$. A coherent $\mathcal{F}$ on $X$ extends to a coherent sheaf
on $\bar X$ \cite[Exercise II.5.15]{hartshorne-ag}. Applying Theorem~\ref{coherent-sheaves:theorem-serre-generation} to the extensions of
$\mathcal{F}, \mathcal{F} \otimes \mathcal{L}, \ldots, \mathcal{F} \otimes \mathcal{L}^{m-1}$ shows that $\mathcal{F} \otimes \mathcal{L}^n$ is globally generated for $n \gg 0$.

Conversely, let $\mathcal{L}$ be ample. \emph{Step 1: affine cover by sets $X_s$.} Let $P \in X$, $U$ an affine neighbourhood on which $\mathcal{L}$ is trivial, and $\mathcal{I}$
the ideal sheaf of the reduced closed subscheme $X \setminus U$. For some $n > 0$, $\mathcal{I} \otimes \mathcal{L}^n$ is globally generated, so there is
$s \in \Gamma(X, \mathcal{I} \otimes \mathcal{L}^n) \subseteq \Gamma(X, \mathcal{L}^n)$ with $s_P \notin \mathfrak{m}_P\mathcal{L}^n_P$. Then $P \in X_s \subseteq U$, and $X_s$ is a basic
open subset of the affine $U$ (where $s$ becomes a function), hence affine. By quasi-compactness finitely many $X_{s_1}, \ldots, X_{s_k}$ cover $X$, and replacing
the $s_i$ by powers we may assume they are all sections of the same $\mathcal{L}^n$.

\emph{Step 2: coordinates.} Let $B_i = \Gamma(X_{s_i}, \mathcal{O})$, a finitely generated $A$-algebra with generators $b_{ij}$. By
Lemma~\ref{coherent-sheaves:lemma-extend-sections}(2) there is $N$ such that $s_i^Nb_{ij}$ extends to a global section $c_{ij}$ of $\mathcal{L}^{nN}$ for all $i$, $j$. The
sections $s_i^N$ and $c_{ij}$ of $\mathcal{L}^{nN}$ generate it, because the $s_i^N$ have no common zero. Let $\varphi \colon X \to \PP^M_A$ be the morphism they define
(Proposition~\ref{divisors:proposition-morphisms-projective}), with homogeneous coordinates $y_i$, $y_{ij}$ corresponding to $s_i^N$, $c_{ij}$.

\emph{Step 3: $\varphi$ is an immersion.} We have $\varphi^{-1}(D_+(y_i)) = X_{s_i}$, and the map $X_{s_i} \to D_+(y_i)$ corresponds to the ring map
$A[y/y_i] \to B_i$ sending $y_{ij}/y_i$ to $c_{ij}/s_i^N = b_{ij}$. It is surjective, so $X_{s_i} \to D_+(y_i)$ is a closed immersion. Hence $\varphi$ is a closed
immersion of $X$ into the open subset $\bigcup_iD_+(y_i)$ of $\PP^M_A$, that is, an immersion, and $\mathcal{L}^{nN} = \varphi^*\mathcal{O}(1)$ is very ample.
\end{proof}

\begin{theorem}[Cohomological criterion]\label{divisors:theorem-cohomological-ampleness}
Let $X$ be proper over a Noetherian ring $A$ and $\mathcal{L}$ invertible. Then $\mathcal{L}$ is ample if and only if for every coherent $\mathcal{F}$ there is $n_0$ with
$H^i(X, \mathcal{F} \otimes \mathcal{L}^n) = 0$ for all $i > 0$ and $n \geq n_0$.
\end{theorem}

\begin{proof}
If $\mathcal{L}$ is ample, some $\mathcal{L}^m$ is very ample (Theorem~\ref{divisors:theorem-ample-very-ample}). As $X$ is proper, the immersion is closed, so $X$ is
projective. Apply Theorem~\ref{coherent-cohomology:theorem-finiteness}(2) to $\mathcal{F} \otimes \mathcal{L}^r$, $0 \leq r < m$.

Conversely, let $P$ be a closed point with ideal sheaf $\mathcal{I}_P$. From $0 \to \mathcal{I}_P\mathcal{F} \to \mathcal{F} \to \mathcal{F} \otimes k(P) \to 0$ and
$H^1(\mathcal{I}_P\mathcal{F} \otimes \mathcal{L}^n) = 0$ for $n$ large, $\Gamma(\mathcal{F} \otimes \mathcal{L}^n) \to \mathcal{F} \otimes \mathcal{L}^n \otimes k(P)$ is onto. By Nakayama's lemma
(Proposition~\ref{coherent-sheaves:proposition-fibre-dimension}(1)), $\mathcal{F} \otimes \mathcal{L}^n$ is then generated by global sections on a neighbourhood of $P$,
a neighbourhood that depends on $n$. Apply this to $\mathcal{F} = \mathcal{O}_X$ to find $n_0$ and a neighbourhood $V$ of $P$ on which $\mathcal{L}^{n_0}$ is globally
generated. Apply it to $\mathcal{F} \otimes \mathcal{L}^r$, $0 \leq r < n_0$, to find $n_1$ and neighbourhoods $V_r$ on which $\mathcal{F} \otimes \mathcal{L}^{r + n_1}$ is globally
generated. On $W_P = V \cap \bigcap_rV_r$, every $\mathcal{F} \otimes \mathcal{L}^n$ with $n \geq n_1$ is globally generated, being a tensor product
$(\mathcal{F} \otimes \mathcal{L}^{r + n_1}) \otimes (\mathcal{L}^{n_0})^{q}$. The $W_P$ for closed points $P$ cover the quasi-compact $X$, since every nonempty closed
subset contains a closed point, so finitely many suffice.
\end{proof}

\begin{example}\label{divisors:example-ample}
\begin{enumerate}
\item $\mathcal{O}(d)$ on $\PP^n_k$, $n \geq 1$, is ample if and only if $d > 0$: for $d > 0$ it is very ample (Example~\ref{varieties:example-projective}(2)), and for
$d \leq 0$ and $\mathcal{F} = \mathcal{O}(-1)$ the sheaf $\mathcal{F} \otimes \mathcal{O}(d)^m = \mathcal{O}(dm - 1)$ is nonzero without nonzero global sections
(Proposition~\ref{coherent-sheaves:proposition-twisting}).
\item The restriction of an ample sheaf to a closed subscheme is ample, and a tensor product of ample sheaves is ample.
\item For a smooth projective variety over $\CC$, a line bundle is ample if and only if the associated holomorphic line bundle is positive
(Definition~\ref{holomorphic-bundles:definition-positive}). This combines Theorem~\ref{kodaira:theorem-embedding} with GAGA (Chapter~\ref{gaga}).
\end{enumerate}
\end{example}

\section{Intersection numbers on surfaces}\label{divisors:section-intersection-surfaces}

\begin{definition}\label{divisors:definition-intersection-number}
Let $X$ be a nonsingular projective surface over an algebraically closed field $k$. For invertible sheaves $\mathcal{L}$, $\mathcal{M}$ put
\[
(\mathcal{L} \cdot \mathcal{M}) = \chi(\mathcal{O}_X) - \chi(\mathcal{L}^{-1}) - \chi(\mathcal{M}^{-1}) + \chi(\mathcal{L}^{-1} \otimes \mathcal{M}^{-1}),
\]
and for divisors $D \cdot E = (\mathcal{O}(D) \cdot \mathcal{O}(E))$, $D^2 = D \cdot D$.
\end{definition}

\begin{proposition}\label{divisors:proposition-intersection-curve}
The pairing $(D, E) \mapsto D \cdot E$ is symmetric, and for an effective divisor $C$ (for instance a curve) and any divisor $D$,
\[
D \cdot C = \chi(\mathcal{O}_C) - \chi(\mathcal{O}_X(-D)|_C) .
\]
\end{proposition}

\begin{proof}
Symmetry is clear. Tensor $0 \to \mathcal{O}(-C) \to \mathcal{O} \to \mathcal{O}_C \to 0$ with $\mathcal{O}$ and with $\mathcal{O}(-D)$ and use additivity of $\chi$:
$\chi(\mathcal{O}) - \chi(\mathcal{O}(-C)) = \chi(\mathcal{O}_C)$ and $\chi(\mathcal{O}(-D)) - \chi(\mathcal{O}(-D-C)) = \chi(\mathcal{O}(-D)|_C)$. Subtracting the second equation from the first
gives the formula.
\end{proof}

\begin{theorem}\label{divisors:theorem-intersection-pairing}
The pairing $(D, E) \mapsto D \cdot E$ is bilinear on $\Pic(X)$. For an irreducible curve $C$, $D \cdot C = \deg(\mathcal{O}(D)|_{\tilde C})$, the degree of the pullback to the
normalisation. If $C$ and $D$ are curves meeting transversally in finitely many points, then $C \cdot D = \#(C \cap D)$, and if $C$, $D$ are distinct irreducible
curves, $C \cdot D \geq 0$.
\end{theorem}

This is \cite[Theorem V.1.1, Proposition V.1.4 and Exercise V.1.1]{hartshorne-ag}; see also \cite[Chapter I]{beauville-surfaces}, where
Definition~\ref{divisors:definition-intersection-number} is taken as the definition. By Proposition~\ref{divisors:proposition-intersection-curve} and the Riemann--Roch
theorem on curves (Chapter~\ref{curves}), $D \cdot C$ is the degree of $\mathcal{O}(D)|_C$, which is additive in $D$. Linearity in the other variable follows by writing any
divisor as a difference of very ample ones, which are linearly equivalent to nonsingular curves by Bertini's theorem.

\begin{example}\label{divisors:example-intersections}
\begin{enumerate}
\item On $\PP^2$, $H \cdot H = \chi(\mathcal{O}) - 2\chi(\mathcal{O}(-1)) + \chi(\mathcal{O}(-2)) = 1 - 0 + 0 = 1$ by Theorem~\ref{coherent-cohomology:theorem-projective-space}. By
bilinearity, curves of degrees $d$ and $e$ satisfy $C \cdot D = de$ (\emph{B\'ezout's theorem}).
\item On $\PP^1 \times \PP^1$, with fibres $F_1 = \{pt\} \times \PP^1$ and $F_2 = \PP^1 \times \{pt\}$, $\chi(\mathcal{O}(a, b)) = (a + 1)(b + 1)$ by the K\"unneth formula for the
\v{C}ech complexes, and the definition gives $F_1^2 = F_2^2 = 0$ and $F_1 \cdot F_2 = 1$.
\item If $\pi \colon \tilde X \to X$ is the blow-up of a point with exceptional curve $E \cong \PP^1$, then $E^2 = -1$, since $\mathcal{O}(E)|_E \cong \mathcal{O}_{\PP^1}(-1)$
(Lemma~\ref{kodaira:lemma-blowup-facts}(1) in the complex case; the algebraic proof is the same).
\end{enumerate}
\end{example}

\begin{theorem}[Hodge index theorem]\label{divisors:theorem-hodge-index}
Let $H$ be an ample divisor on a nonsingular projective surface $X$. If $D \cdot H = 0$ and $D$ is not numerically trivial (Definition~\ref{divisors:definition-num}),
then $D^2 < 0$. Equivalently, the intersection form on $\mathrm{Num}(X) \otimes \RR$ has signature $(1, \rho - 1)$.
\end{theorem}

This is \cite[Theorem V.1.9]{hartshorne-ag}, proved there with Riemann--Roch on surfaces. Over $\CC$ it is the part of
Corollary~\ref{lefschetz-theorems:corollary-surfaces} concerning classes of divisors, once intersection numbers are identified with cup products of cycle classes
(Chapter~\ref{cycle-class}).

\begin{theorem}[Nakai--Moishezon criterion]\label{divisors:theorem-nakai}
A divisor $D$ on a nonsingular projective surface is ample if and only if $D^2 > 0$ and $D \cdot C > 0$ for every irreducible curve $C$.
\end{theorem}

This is \cite[Theorem V.1.10]{hartshorne-ag}. In higher dimensions the criterion reads $D^{\dim V} \cdot V > 0$ for every closed subvariety $V$
\cite[Theorem 1.2.23]{lazarsfeld-positivity-1}.

\section{The N\'eron--Severi group}\label{divisors:section-neron-severi}

\begin{definition}\label{divisors:definition-num}
Let $X$ be a nonsingular projective variety over an algebraically closed field. A divisor $D$ is \emph{numerically trivial}, $D \equiv 0$, if $\deg(\mathcal{O}(D)|_C) = 0$ for every
irreducible curve $C \subseteq X$. $\mathrm{Num}(X) = \Pic(X)/\equiv$. Two divisors are \emph{algebraically equivalent} if they are members of a family of divisors
parametrised by a connected base, and $\mathrm{NS}(X) = \Pic(X)/\{\text{algebraic equivalence}\}$ is the \emph{N\'eron--Severi group}. The rank of $\mathrm{Num}(X)$
is the \emph{Picard number} $\rho(X)$.
\end{definition}

\begin{theorem}[Theorem of the base]\label{divisors:theorem-base}
$\mathrm{NS}(X)$ is a finitely generated abelian group, algebraic equivalence implies numerical equivalence, and $\mathrm{Num}(X)$ is a free abelian group of finite rank
$\rho(X)$, the quotient of $\mathrm{NS}(X)$ by its torsion.
\end{theorem}

This is due to Severi and N\'eron; see \cite[Section 1.1]{lazarsfeld-positivity-1}. Over $\CC$ it follows from Hodge theory. The kernel $\Pic^0(X)$ of the first Chern
class is the Picard torus (Remark~\ref{holomorphic-bundles:remark-picard-variety}), and $\mathrm{NS}(X)$ embeds into $H^2(X; \ZZ)$ with image the integral classes of type
$(1,1)$ (Chapter~\ref{lefschetz-11}). So $\rho(X) \leq h^{1,1}(X)$.

\begin{example}\label{divisors:example-picard-numbers}
\begin{enumerate}
\item $\rho(\PP^n) = 1$ and $\rho(\PP^1 \times \PP^1) = 2$, since $\Pic(\PP^1 \times \PP^1) \cong \ZZ^2$, generated by the two rulings.
\item For a quartic surface $X \subseteq \PP^3$, $1 \leq \rho(X) \leq h^{1,1} = 20$ (Example~\ref{hodge-decomposition:example-quartic}). The very general quartic has $\rho = 1$
(Noether--Lefschetz), and the Fermat quartic $x^4 + y^4 + z^4 + w^4 = 0$ has $\rho = 20$.
\item For $E \times E'$ a product of elliptic curves, $\rho = 2 + \operatorname{rank}\Hom(E, E')$, which is $2$, $3$ or $4$ (Chapter~\ref{abelian-varieties}).
\end{enumerate}
\end{example}

\section{Exercises}\label{divisors:section-exercises}

\begin{exercise}\label{divisors:exercise-cl-minus-hypersurface}
Let $Y \subseteq \PP^n_k$ be an irreducible hypersurface of degree $d$. Show that $\mathrm{Cl}(\PP^n \setminus Y) \cong \ZZ/d\ZZ$.
\end{exercise}

\begin{solution}
By Proposition~\ref{divisors:proposition-class-group}(2), $\mathrm{Cl}(\PP^n \setminus Y) = \mathrm{Cl}(\PP^n)/\ZZ[Y] = \ZZ/d\ZZ$, using (3) and $\deg Y = d$.
\end{solution}

\begin{exercise}\label{divisors:exercise-linear-system-conics}
Show that the conics in $\PP^2$ form a linear system $|2H|$ of dimension $5$ without base points, and that the morphism it defines is the Veronese embedding
$v_2 \colon \PP^2 \to \PP^5$.
\end{exercise}

\begin{solution}
$\Gamma(\PP^2, \mathcal{O}(2))$ is the $6$-dimensional space of quadratic forms (Proposition~\ref{coherent-sheaves:proposition-twisting}), so $\dim|2H| = 5$ by
Proposition~\ref{divisors:proposition-linear-system}. The monomials $x_i^2$ have no common zero, so there are no base points, and the morphism of
Proposition~\ref{divisors:proposition-morphisms-projective} given by the six quadratic monomials is $v_2$.
\end{solution}

\begin{exercise}\label{divisors:exercise-negative-curve}
Let $C$ be an irreducible curve on a nonsingular projective surface with $C^2 < 0$. Show that $C$ is the only member of $|C|$, that is, $h^0(\mathcal{O}(C)) = 1$.
\end{exercise}

\begin{solution}
Let $D \in |C|$ and write $D = aC + D'$ with $a \geq 0$ and $D'$ effective not containing $C$. If $a \geq 1$, then $(a - 1)C + D'$ is effective and linearly
equivalent to $0$. By Proposition~\ref{divisors:proposition-linear-system} with $h^0(\mathcal{O}) = 1$, the only effective divisor in $|0|$ is $0$, so $a = 1$, $D' = 0$ and
$D = C$. If $a = 0$, then $C^2 = C \cdot D' \geq 0$ by Theorem~\ref{divisors:theorem-intersection-pairing}, since $D'$ is a sum of curves different from $C$. This
contradicts $C^2 < 0$.
\end{solution}

\begin{exercise}\label{divisors:exercise-ample-restriction}
Show that if $\mathcal{L}$ is ample on $X$ and $Y \subseteq X$ is a closed subscheme, then $\mathcal{L}|_Y$ is ample, using Theorem~\ref{divisors:theorem-cohomological-ampleness}
(for $X$ proper over a Noetherian ring).
\end{exercise}

\begin{solution}
For coherent $\mathcal{G}$ on $Y$, $i_*\mathcal{G}$ is coherent on $X$ and $H^j(Y, \mathcal{G} \otimes \mathcal{L}|_Y^n) = H^j(X, i_*\mathcal{G} \otimes \mathcal{L}^n)$ by the projection formula and
Example~\ref{sheaf-cohomology:example-leray}. These vanish for $j > 0$ and $n \gg 0$.
\end{solution}
